\FloatBarrier
\section{FT outer mount}
\label{sec:outer-mount}

The eighth beamline study renames outerMount to FTOuterMount across ten source and runtime CAD variations
and ten GEMC2 SQLite rows. The stainless-steel material, color, mother, and passive placement are retained.
The mount remains at z = 1273.27 mm and the unchanged FTOuterFlange at z = 1272.77 mm, both with the
original passive 180-degree y rotation. Their 0.5 mm relative translation is included in every clearance
calculation.

The original mount contains 3,094 triangles, twelve boundary edges forming four triangular surface gaps, and
one manifold surface component. Its open surface has no defined solid volume. Closing these four missing
faces gives a 3,098-triangle sealed reference of genus 11 and volume 3889831.731170 mm$^3$. The central
passage and existing drilled openings are retained. This repaired reference defines the reported
material-volume comparison; the original engineering solid mass is not independently certified.

The original mount/flange check reports small interference at their nominal 95.25 mm radius because the two
cylindrical surfaces have different polygonal sampling. The flange inner polygon has minimum inradius
95.200045 mm. The mount lower outer rim is reduced to 95.150045 mm, approximately 0.100 mm of radial relief.
Adjacent sidewall facets interpolate toward unchanged upper vertices. All 883 facets adjoining the relieved
end are protected at their saved float32 coordinates during topology-preserving quadric simplification. The
flange STL is unchanged, and no placement is shifted to conceal the interference.

The final mount contains 2,740 triangles, 11.44\% fewer than the raw engineering mesh, and 137,084 bytes
rather than 154,784. It is one closed manifold component of genus 11 with no detected self-intersections or
zero-area triangles. Bidirectional sampling against the relieved reference gives a maximum difference of
0.042078 mm, below the 0.05 mm guard. These are sampled differences rather than proven global error bounds.
Volume changes by -0.249235\% against the sealed reference; simplification alone changes the relieved
reference volume by 0.00011056\%. The axial bounds are retained, with at most about 0.00065 mm change in the
transverse bounding coordinates. The builder reproduces the final STL byte for byte from the archived
original.

Table~\ref{tab:outer-mount-mesh} records the original and optimized mount inventory.

\begin{table}[htbp]
\centering
\caption{Original and optimized FT outer mount.}
\label{tab:outer-mount-mesh}
\begin{tabular}{lrr}
\toprule
Quantity & Original & Optimized \\
\midrule
Triangles & 3,094 & 2,740 \\
Bytes & 154,784 & 137,084 \\
Boundary edges & 12 & 0 \\
Surface components & 1 & 1 \\
Solid volume (mm$^3$) & Undefined & 3880136.906 \\
\bottomrule
\end{tabular}
\end{table}

Placed-mesh verification finds zero triangle contacts and no contained vertices in either direction between
mount and flange. In mount coordinates, the flange is shifted by +0.5 mm along z. Each mount triangle is
clipped to the flange z interval, and convexity of the radial norm bounds its radius by the clipped polygon
vertices. The maximum is 95.164071 mm, giving at least 0.035974 mm radial separation from the flange inner
polygon throughout their common axial interval. The minimum bidirectional vertex-to-surface distance is
0.085929 mm; this latter value is sampled and is not presented as a global surface-clearance bound. All ten
placement variants retain the same verified relative transform.

A twelve-volume assembly fixture includes the mount, unchanged inner shield and flange, and nine copper
bolts. With both original and optimized fixtures using the new mount name and identical placements, a
10,000-point Geant4 test gives five flange/mount overlap reports and one surface-hole warning for the
original mount. The optimized assembly completes with no Geant4 exceptions or overlap reports, resolving the
joint interference documented in the preceding flange study. Full default beamline construction also
completes without missing meshes or a mount/flange overlap. Its three solid warnings and eight overlap
reports concern other geometry, so this study does not establish complete beamline overlap freedom. All
fifteen native geometry/material variations and fourteen CAD variations pass the semantic comparison against
GEMC2.

The tracking benchmark isolates this one mount in air. Both cases use the new name, the same material and
placement, and the unchanged fixed GEMC3 executable snapshot from the flange study. Matched 5 GeV LUND
primaries begin at the target and sample solid angle uniformly over 1.5--2.4 degrees and full azimuth.
FTFP\_BERT runs without optical processes or magnetic field, using one thread, default production cuts, and
no GUI or event output. Five alternating-order repetitions use seeds 12345--12349 and 200 events per
electron, proton, positive-pion, and neutron run. Elapsed time since beamOn measures transport; separate
zero-event samples measure startup and peak RSS.

Table~\ref{tab:outer-mount-timing} records median transport wall seconds per 200 events.

\begin{table}[htbp]
\centering
\caption{Median transport wall seconds over five repetitions.}
\label{tab:outer-mount-timing}
\begin{tabular}{lrrr}
\toprule
Particle & Original & Optimized & Original/optimized \\
\midrule
e- & 15.97580 & 15.51250 & 1.030 \\
proton & 2.13591 & 2.10357 & 1.015 \\
pi+ & 2.93950 & 2.85889 & 1.028 \\
neutron & 1.83088 & 1.93090 & 0.948 \\
\bottomrule
\end{tabular}
\end{table}

Median original/optimized transport wall-time ratios are 1.030 for electrons, 1.015 for protons, 1.028 for
positive pions, and 0.948 for neutrons. A ratio below one denotes longer optimized tracking time. All 52
processes complete normally, and all forty transport runs report the requested 200 events. Startup
simulation wall time changes from 0.705080 to 0.693238 s, and peak startup RSS changes from 120.12 to 119.38
MiB. Closure repairs and relief change material boundaries, so these short workload timings are not
attributed solely to triangle reduction. They do not establish detector-response equivalence or full CLAS12
speedup. The preceding flange control at fixed material volume remains separately reported.

Table~\ref{tab:outer-mount-warnings} gives isolated-fixture exceptions across all 52 processes.

\begin{table}[htbp]
\centering
\caption{Geant4 exception counts in the isolated mount fixtures.}
\label{tab:outer-mount-warnings}
\begin{tabular}{lrr}
\toprule
Diagnostic & Original & Optimized \\
\midrule
GeomSolids1001 & 26 & 0 \\
GeomNav1002 & 0 & 0 \\
GeomNav0003 & 0 & 0 \\
GeomVol1002 & 0 & 0 \\
\bottomrule
\end{tabular}
\end{table}

\FloatBarrier
Figure~\ref{fig:outer-mount-stl} compares the stored mount triangulations.

\begin{figure}[htbp]
\centering
\includegraphics[width=\textwidth]{figures/outer_mount_stl_comparison.pdf}
\caption{Original FT outer mount (left) and repaired, relieved, and simplified mount (right).
Matched cameras and bounds use original STL coordinates; existing drilled openings are retained.}
\label{fig:outer-mount-stl}
\end{figure}
